% =====================================================================
%  Relatório MODSI — Sistema de Gestão de Festivais
%  Compilar com pdfLaTeX (ex.: Overleaf). Ficheiro principal: main.tex
% =====================================================================
\documentclass[12pt,a4paper]{report}

\usepackage[utf8]{inputenc}
\usepackage[T1]{fontenc}
\usepackage[portuguese]{babel}
\usepackage[a4paper,margin=2.5cm]{geometry}
\usepackage{graphicx}
\usepackage{xcolor}
\usepackage{array}
\usepackage{tabularx}
\usepackage{longtable}
\usepackage{float}
\usepackage{caption}
\usepackage{fancyhdr}
\setlength{\headheight}{15pt}
\usepackage{tikz}
\usetikzlibrary{shapes.geometric,arrows.meta,positioning,calc,babel}
\usepackage[hidelinks]{hyperref}

\addto\captionsportuguese{\renewcommand{\contentsname}{Índice}}

% ---------- Cabeçalho ----------
\pagestyle{fancy}
\fancyhf{}
\fancyhead[L]{\small\itshape\leftmark}
\fancyfoot[C]{\thepage}
\renewcommand{\chaptermark}[1]{\markboth{\thechapter. #1}{}}

% ---------- Nota de trabalho por fazer (remover antes da entrega) ----------
\newcommand{\pendente}[1]{%
  \par\noindent\fcolorbox{red!70!black}{red!5}{%
  \parbox{\dimexpr\linewidth-2\fboxsep-2\fboxrule}{\small\color{red!60!black}%
  \textbf{A completar:} #1}}\par\medskip}

% ---------- Cenários de casos de uso ----------
% \cenario{UCxx}{Nome}{Ator}{Pré-condição}{Fluxo principal}{Pós-condição}{Alternativos}
\newcommand{\cenario}[7]{%
  \par\noindent
  \begin{tabularx}{\textwidth}{|>{\bfseries}p{3.1cm}|X|}
    \hline
    \multicolumn{2}{|l|}{\textbf{#1 -- #2}}\\ \hline
    Ator            & #3 \\ \hline
    Pré-condição    & #4 \\ \hline
    Fluxo principal & #5 \\ \hline
    Pós-condição    & #6 \\ \hline
    Alternativos    & #7 \\ \hline
  \end{tabularx}\par\medskip}

% ---------- Estilos TikZ ----------
\tikzset{
  uc/.style={ellipse, draw, fill=blue!8, minimum width=5.6cm,
             minimum height=0.85cm, font=\footnotesize, align=center, inner sep=1pt},
  sistema/.style={draw, thick, rounded corners=2pt},
  assoc/.style={thin},
  relacao/.style={dashed, -{Stealth[length=2mm]}, thin},
  relnome/.style={font=\scriptsize, right, align=left},
  general/.style={-{Triangle[open,length=3mm,width=3mm]}, thin},
  estado/.style={draw, rounded corners=6pt, fill=blue!8, minimum width=2.6cm,
                 minimum height=0.9cm, font=\small},
  trans/.style={-{Stealth[length=2mm]}, thin},
  translab/.style={font=\scriptsize, align=center},
}
\newcommand{\incl}{\guillemotleft include\guillemotright}
\newcommand{\extd}{\guillemotleft extend\guillemotright}

% Ator UML: \ator{nome-coordenada}{posição}{rótulo}
\newcommand{\ator}[3]{%
  \begin{scope}[shift={(#2)}]
    \draw[thick] (0,0.62) circle (0.17);
    \draw[thick] (0,0.45) -- (0,-0.1);
    \draw[thick] (-0.32,0.28) -- (0.32,0.28);
    \draw[thick] (0,-0.1) -- (-0.27,-0.55);
    \draw[thick] (0,-0.1) -- (0.27,-0.55);
    \node[below, font=\small, align=center] at (0,-0.55) {#3};
    \coordinate (#1) at (0.35,0.15);
    \coordinate (#1-left) at (-0.35,0.15);
    \coordinate (#1-top) at (0,0.8);
    \coordinate (#1-bottom) at (0,-1.0);
  \end{scope}}

\begin{document}

% =====================================================================
\begin{titlepage}
  \centering
  \vspace*{3cm}
  {\Large Modelação e Simulação Industrial\par}
  \vspace{4cm}
  \rule{\textwidth}{0.4pt}\par\vspace{0.6cm}
  {\LARGE Sistema de Gestão de Festivais\par}
  \vspace{0.6cm}\rule{\textwidth}{0.4pt}\par
  \vspace{1.5cm}
  {\large Grupo \textbf{XX}\par\vspace{0.3cm}
   Nome Apelido -- 12XXXXX\\
   Nome Apelido -- 12XXXXX\\
   Nome Apelido -- 12XXXXX\par}
  \vspace{1.0cm}
  {\large Docente: Nome do Docente\par}
  \vfill
  \IfFileExists{isep_logo.png}{\includegraphics[width=6cm]{isep_logo.png}\par\vspace{1cm}}{}
  Mestrado em Engenharia Eletrotécnica e de Computadores\\[0.2cm]
  Departamento de Engenharia Eletrotécnica\\[0.6cm]
  Ano letivo 2026/2027
\end{titlepage}

\pagenumbering{roman}
\tableofcontents
\listoffigures
\listoftables
\clearpage
\pagenumbering{arabic}

\chapter{Introdução}

\section{Contextualização}

A organização de festivais envolve a coordenação de um elevado volume de
informação: datas e locais, palcos, artistas e respetiva programação,
inscrições, pagamentos e o controlo de acesso ao recinto. Quando esta
informação é gerida de forma dispersa (folhas de cálculo, mensagens, listas em
papel), surgem erros frequentes, como inscrições duplicadas, bilhetes
reutilizados ou programações com conflitos de horário.

O presente trabalho propõe um \textbf{Sistema de Gestão de Festivais}: uma
aplicação \emph{web} que centraliza esta informação numa base de dados
relacional e disponibiliza, a cada tipo de utilizador, apenas as
funcionalidades que lhe dizem respeito. O público consulta festivais e
programação; os participantes inscrevem-se, efetuam um pagamento simulado e
recebem um bilhete com código QR; os organizadores gerem os seus festivais; o
\emph{Staff} valida as entradas no recinto; e o administrador supervisiona a
plataforma.

\section{Objetivos}

Os objetivos do trabalho são:
\begin{itemize}
  \item modelar o sistema de informação, identificando atores, requisitos,
        casos de uso e respetivos cenários;
  \item conceber e implementar uma base de dados relacional que garanta a
        integridade dos dados (inscrições únicas, bilhetes de utilização
        única, programação sem conflitos);
  \item desenhar a interface gráfica (\emph{wireframes}, formatos de saída e
        \emph{wireflow});
  \item implementar uma aplicação \emph{web} responsiva e multilingue
        (Português e Inglês), com autenticação local e controlo de acesso por
        perfil;
  \item gerir o projeto de forma colaborativa, com controlo de versões e um
        quadro de tarefas.
\end{itemize}

\section{Âmbito}

Ficam dentro do âmbito: a gestão de festivais, palcos, artistas e programação;
o registo e autenticação de utilizadores; as inscrições com pagamento
\textbf{simulado}; a emissão de bilhetes com código QR; a validação de entradas;
e a administração de contas e perfis.

Ficam fora do âmbito: pagamentos reais e reembolsos, venda de diferentes tipos
de bilhete (por exemplo, diário ou VIP) e integração com sistemas externos de
bilhética.

\section{Metodologia e Gestão do Projeto}

\pendente{Descrever a abordagem iterativa adotada (Scrum/Kanban), o
\emph{Product Backlog} com as histórias de utilizador e critérios de aceitação,
a distribuição de tarefas pelos elementos do grupo e as ferramentas usadas
(GitHub, Jira ou quadro Kanban). Incluir uma captura do quadro.}

\chapter{Análise de Requisitos}

\section{Definição de Atores}

O sistema distingue cinco atores. Os quatro atores autenticados são
especializações de um ator abstrato, o \emph{Utilizador Autenticado}, que
reúne as funcionalidades comuns (iniciar e terminar sessão e editar o perfil).
Cada conta tem um único perfil. O registo público cria sempre uma conta com
perfil \emph{Participante}; os perfis \emph{Organizador}, \emph{Staff} e
\emph{Administrador} são atribuídos pelo Administrador.

\begin{table}[H]
\centering
\caption{Atores do sistema}
\label{tab:atores}
\begin{tabularx}{\textwidth}{|l|X|}
\hline
\textbf{Ator} & \textbf{Descrição} \\ \hline
Visitante & Utilizador não autenticado. Consulta a informação pública
  (festivais publicados, programação, palcos e artistas) e pode registar uma
  conta. \\ \hline
\emph{Utilizador Autenticado} & Ator abstrato que generaliza os quatro perfis
  seguintes. Inicia e termina sessão e edita os seus dados, palavra-passe e
  idioma. \\ \hline
Participante & Utilizador registado que se inscreve em festivais, efetua o
  pagamento simulado, consulta as suas inscrições e bilhetes QR e cancela
  inscrições dentro do prazo. \\ \hline
Organizador & Responsável pelos festivais que cria: dados do festival, palcos,
  artistas, programação, associação de membros do Staff e acompanhamento das
  inscrições. \\ \hline
Staff & Colaborador associado a um ou mais festivais pelo Organizador.
  Consulta os participantes, verifica bilhetes e regista as entradas no
  recinto. \\ \hline
Administrador & Responsável pela gestão global da plataforma: contas,
  perfis, ativação de festivais e indicadores globais. \\ \hline
\end{tabularx}
\end{table}

\section{Requisitos Funcionais}

A Tabela~\ref{tab:rf} lista os requisitos funcionais. A coluna \emph{UC}
indica os casos de uso (Capítulo~\ref{cap:modelacao}) que concretizam cada
requisito. A prioridade segue a classificação \emph{Alta} (indispensável para
a demonstração) e \emph{Média} (a implementar se o tempo o permitir).

\begin{longtable}{|l|p{8.6cm}|l|l|}
\caption{Requisitos funcionais}\label{tab:rf}\\
\hline
\textbf{ID} & \textbf{O sistema deve permitir\dots} & \textbf{Prior.} & \textbf{UC} \\ \hline
\endfirsthead
\hline
\textbf{ID} & \textbf{O sistema deve permitir\dots (cont.)} & \textbf{Prior.} & \textbf{UC} \\ \hline
\endhead
\multicolumn{4}{c}{\textbf{Visitante}} \\ \hline
RF01 & listar os festivais publicados e ativos. & Alta & UC01 \\ \hline
RF02 & pesquisar festivais pelo nome e filtrá-los por data e localização. & Média & UC02, UC03 \\ \hline
RF03 & consultar os detalhes de um festival, incluindo programação, palcos e artistas. & Alta & UC04, UC05 \\ \hline
RF04 & registar uma conta, que fica com o perfil Participante. & Alta & UC06 \\ \hline
\multicolumn{4}{c}{\textbf{Utilizador Autenticado}} \\ \hline
RF05 & iniciar e terminar sessão com email e palavra-passe (autenticação local). & Alta & UC07, UC08 \\ \hline
RF06 & editar os dados pessoais, a palavra-passe e o idioma preferido. & Média & UC09 \\ \hline
\multicolumn{4}{c}{\textbf{Participante}} \\ \hline
RF07 & inscrever-se num festival, respeitando o período de inscrições e a lotação. & Alta & UC10 \\ \hline
RF08 & efetuar um pagamento simulado, registando cada tentativa e o seu resultado. & Alta & UC11 \\ \hline
RF09 & emitir automaticamente um bilhete com código QR único quando o pagamento é aprovado. & Alta & UC11, UC13 \\ \hline
RF10 & consultar as suas inscrições e respetivos estados. & Alta & UC12 \\ \hline
RF11 & cancelar uma inscrição até à data limite de cancelamento, revogando o bilhete. & Média & UC14 \\ \hline
\multicolumn{4}{c}{\textbf{Organizador}} \\ \hline
RF12 & criar e editar festivais (dados, datas, período de inscrições, capacidade e preço). & Alta & UC15, UC16 \\ \hline
RF13 & alterar o estado de um festival (publicar, cancelar, terminar). & Alta & UC16 \\ \hline
RF14 & gerir os palcos de um festival. & Alta & UC17 \\ \hline
RF15 & gerir os artistas e a sua associação ao festival. & Alta & UC18 \\ \hline
RF16 & gerir a programação (atuações), impedindo conflitos de horário no mesmo palco. & Alta & UC19 \\ \hline
RF17 & associar e remover membros do Staff de um festival. & Alta & UC20 \\ \hline
RF18 & consultar as inscrições do festival e respetivos indicadores (gráfico). & Média & UC21 \\ \hline
\multicolumn{4}{c}{\textbf{Staff}} \\ \hline
RF19 & consultar os festivais aos quais está associado. & Alta & UC22 \\ \hline
RF20 & consultar e pesquisar os participantes inscritos num festival. & Alta & UC23 \\ \hline
RF21 & verificar a validade de um bilhete através do código QR ou do código manual. & Alta & UC24 \\ \hline
RF22 & registar a entrada de um participante, impedindo a reutilização do bilhete. & Alta & UC25 \\ \hline
RF23 & consultar o histórico de entradas de um festival. & Média & UC26 \\ \hline
\multicolumn{4}{c}{\textbf{Administrador}} \\ \hline
RF24 & gerir contas de utilizador (criar, editar, ativar e desativar). & Alta & UC27 \\ \hline
RF25 & atribuir e alterar o perfil de um utilizador. & Alta & UC28 \\ \hline
RF26 & consultar todos os festivais e ativá-los ou desativá-los. & Alta & UC29, UC30 \\ \hline
RF27 & consultar indicadores globais de inscrições e pagamentos (\emph{dashboard}). & Média & UC31 \\ \hline
\end{longtable}

\section{Requisitos Não Funcionais}

Os requisitos não funcionais incluem os requisitos de implementação definidos
no enunciado do trabalho.

\begin{longtable}{|l|p{12.4cm}|}
\caption{Requisitos não funcionais}\label{tab:rnf}\\
\hline
\textbf{ID} & \textbf{Requisito não funcional} \\ \hline
\endfirsthead
\hline
\textbf{ID} & \textbf{Requisito não funcional (cont.)} \\ \hline
\endhead
RNF01 & \textbf{Segurança:} autenticação local com sessão para todas as funcionalidades reservadas. \\ \hline
RNF02 & \textbf{Segurança:} controlo de acesso por perfil e, para Organizador e Staff, pela associação ao festival. \\ \hline
RNF03 & \textbf{Segurança:} palavras-passe guardadas apenas como \emph{hash} (\texttt{password\_hash}/bcrypt); consultas SQL parametrizadas (\emph{prepared statements}). \\ \hline
RNF04 & \textbf{Arquitetura:} código organizado segundo o padrão MVC (\emph{Model-View-Controller}). \\ \hline
RNF05 & \textbf{Dados:} base de dados relacional MySQL acedida por PDO, com chaves estrangeiras, restrições de unicidade, vistas e transações. \\ \hline
RNF06 & \textbf{Interatividade:} pesquisa, filtros e validação de bilhetes efetuados por AJAX, sem recarregar a página. \\ \hline
RNF07 & \textbf{Usabilidade:} interface responsiva (computador, \emph{tablet} e telemóvel) com Bootstrap. \\ \hline
RNF08 & \textbf{Usabilidade:} interface multilingue, pelo menos em Português e Inglês. \\ \hline
RNF09 & \textbf{Usabilidade:} mensagens claras de sucesso e de erro em todas as operações. \\ \hline
RNF10 & \textbf{Integridade:} impedir inscrições duplicadas e a reutilização de um bilhete, ao nível da base de dados. \\ \hline
RNF11 & \textbf{Integridade:} pagamentos simulados, sem movimentação de dinheiro real. \\ \hline
RNF12 & \textbf{Desempenho:} listagens com paginação e índices nas colunas de pesquisa. \\ \hline
RNF13 & \textbf{Compatibilidade:} funcionamento nos navegadores modernos (Chrome, Firefox, Edge, Safari). \\ \hline
RNF14 & \textbf{Gestão:} controlo de versões com Git/GitHub e acompanhamento das tarefas num quadro Kanban/Jira. \\ \hline
\end{longtable}

\section{Regras de Negócio}
\label{sec:regras}

As regras seguintes clarificam o comportamento do sistema e são a referência
comum para os cenários (Capítulo~\ref{cap:modelacao}) e para o modelo de dados
(Capítulo~\ref{cap:bd}).

\begin{longtable}{|l|p{12.4cm}|}
\caption{Regras de negócio}\label{tab:rn}\\
\hline
\textbf{ID} & \textbf{Regra} \\ \hline
\endfirsthead
\hline
\textbf{ID} & \textbf{Regra (cont.)} \\ \hline
\endhead
RN01 & Cada utilizador tem um único perfil. O registo público cria contas com perfil Participante; os restantes perfis são atribuídos pelo Administrador. \\ \hline
RN02 & Uma conta desativada não pode iniciar sessão. \\ \hline
RN03 & Um festival só é visível ao público se estiver no estado \emph{publicado} e ativo. O \textbf{estado} é gerido pelo Organizador (rascunho, publicado, cancelado, terminado); o indicador \textbf{ativo} é gerido pelo Administrador. \\ \hline
RN04 & O Organizador só gere os festivais que criou. \\ \hline
RN05 & Só é possível inscrever-se num festival publicado e ativo, entre as datas de abertura e de fecho das inscrições, e enquanto o número de inscrições confirmadas for inferior à capacidade. \\ \hline
RN06 & Cada participante tem no máximo uma inscrição \emph{não cancelada} por festival. Depois de cancelar, pode inscrever-se de novo, o que cria uma nova inscrição; a cancelada fica como histórico, com o respetivo pagamento e bilhete revogado. \\ \hline
RN07 & O valor da inscrição é o preço do festival no momento da inscrição; alterações posteriores do preço não afetam inscrições existentes. \\ \hline
RN08 & Uma inscrição pendente pode ter várias tentativas de pagamento, mas no máximo uma aprovada. A aprovação confirma a inscrição e emite o bilhete numa única transação. \\ \hline
RN09 & Cada inscrição confirmada tem exatamente um bilhete, com um código QR único e não previsível. \\ \hline
RN10 & O cancelamento só é permitido até à data limite de cancelamento e se ainda não existir entrada registada. O bilhete é revogado. Não existe reembolso (pagamentos simulados). \\ \hline
RN11 & Um bilhete é válido para entrada se não estiver revogado, a inscrição estiver confirmada, o festival decorrer nessa data e o bilhete ainda não tiver sido usado. Cada bilhete permite uma única entrada. \\ \hline
RN12 & O Staff só consulta e valida bilhetes dos festivais a que está associado com associação ativa. Remover um membro do Staff desativa a associação, mantendo o histórico. \\ \hline
RN13 & Uma atuação pertence a um artista associado ao festival e a um palco do mesmo festival, decorre dentro das datas do festival e não se sobrepõe a outra atuação no mesmo palco. \\ \hline
RN14 & Não é possível remover um palco ou desassociar um artista que tenham atuações marcadas. \\ \hline
\end{longtable}

\chapter{Modelação do Sistema}
\label{cap:modelacao}

\section{Tabela de Casos de Uso}

A Tabela~\ref{tab:uc} identifica os 31 casos de uso, o ator principal, os
requisitos funcionais que concretizam e as relações
\incl{} e \extd{} entre casos de uso. Cada ator concreto tem pelo menos cinco
casos de uso próprios. Os casos UC07 a UC09 são comuns a todos os
utilizadores autenticados e são herdados através da generalização.

\begin{longtable}{|l|p{5.6cm}|l|l|p{2.6cm}|}
\caption{Casos de uso do sistema}\label{tab:uc}\\
\hline
\textbf{UC} & \textbf{Caso de uso} & \textbf{Ator} & \textbf{RF} & \textbf{Relações} \\ \hline
\endfirsthead
\hline
\textbf{UC} & \textbf{Caso de uso (cont.)} & \textbf{Ator} & \textbf{RF} & \textbf{Relações} \\ \hline
\endhead
UC01 & Consultar lista de festivais & Visitante & RF01 & -- \\ \hline
UC02 & Pesquisar festivais pelo nome & Visitante & RF02 & \extd{} UC01 \\ \hline
UC03 & Filtrar festivais por data ou localização & Visitante & RF02 & \extd{} UC01 \\ \hline
UC04 & Consultar detalhes de um festival & Visitante & RF03 & -- \\ \hline
UC05 & Consultar programação, palcos e artistas & Visitante & RF03 & \extd{} UC04 \\ \hline
UC06 & Registar conta & Visitante & RF04 & -- \\ \hline
UC07 & Iniciar sessão & Util. Autent. & RF05 & -- \\ \hline
UC08 & Terminar sessão & Util. Autent. & RF05 & -- \\ \hline
UC09 & Editar perfil (dados, palavra-passe, idioma) & Util. Autent. & RF06 & -- \\ \hline
UC10 & Inscrever-se num festival & Participante & RF07 & -- \\ \hline
UC11 & Efetuar pagamento simulado & Participante & RF08, RF09 & \extd{} UC10, UC12 \\ \hline
UC12 & Consultar as minhas inscrições & Participante & RF10 & -- \\ \hline
UC13 & Consultar bilhete QR & Participante & RF09 & \extd{} UC12 \\ \hline
UC14 & Cancelar inscrição & Participante & RF11 & \extd{} UC12 \\ \hline
UC15 & Criar festival & Organizador & RF12 & -- \\ \hline
UC16 & Editar festival e alterar estado & Organizador & RF12, RF13 & -- \\ \hline
UC17 & Gerir palcos & Organizador & RF14 & -- \\ \hline
UC18 & Gerir artistas do festival & Organizador & RF15 & -- \\ \hline
UC19 & Gerir programação (atuações) & Organizador & RF16 & -- \\ \hline
UC20 & Gerir Staff do festival & Organizador & RF17 & -- \\ \hline
UC21 & Consultar inscrições do festival & Organizador & RF18 & -- \\ \hline
UC22 & Consultar festivais associados & Staff & RF19 & -- \\ \hline
UC23 & Consultar participantes inscritos & Staff & RF20 & -- \\ \hline
UC24 & Verificar bilhete & Staff & RF21 & -- \\ \hline
UC25 & Validar entrada & Staff & RF22 & \incl{} UC24 \\ \hline
UC26 & Consultar histórico de entradas & Staff & RF23 & -- \\ \hline
UC27 & Gerir contas de utilizador & Administrador & RF24 & -- \\ \hline
UC28 & Atribuir perfil & Administrador & RF25 & \extd{} UC27 \\ \hline
UC29 & Consultar todos os festivais & Administrador & RF26 & -- \\ \hline
UC30 & Ativar ou desativar festival & Administrador & RF26 & \extd{} UC29 \\ \hline
UC31 & Consultar indicadores globais & Administrador & RF27 & -- \\ \hline
\end{longtable}

% =====================================================================
\section{Diagramas de Casos de Uso}

Para manter a legibilidade, o diagrama de casos de uso é apresentado por ator.
A Figura~\ref{fig:uc-geral} mostra a generalização dos atores autenticados e os
casos de uso comuns. As Figuras~\ref{fig:uc-visitante} a~\ref{fig:uc-admin}
mostram os casos de uso de cada ator.

\begin{figure}[H]
\centering
\resizebox{0.9\textwidth}{!}{%
\begin{tikzpicture}
  \ator{par}{0,5.7}{Participante}
  \ator{org}{0,3.8}{Organizador}
  \ator{stf}{0,1.9}{Staff}
  \ator{adm}{0,0}{Administrador}
  \ator{ua}{3.6,2.85}{\emph{Utilizador}\\\emph{Autenticado}}
  \draw[general] (par) -- (ua-left);
  \draw[general] (org) -- (ua-left);
  \draw[general] (stf) -- (ua-left);
  \draw[general] (adm) -- (ua-left);
  \draw[sistema] (5.6,-0.6) rectangle (12.4,5.6);
  \node[font=\small\bfseries] at (9,5.2) {Sistema de Gestão de Festivais};
  \node[uc] (uc07) at (9,4.0) {UC07 Iniciar sessão};
  \node[uc] (uc08) at (9,2.6) {UC08 Terminar sessão};
  \node[uc] (uc09) at (9,1.2) {UC09 Editar perfil};
  \draw[assoc] (ua) -- (uc07.west);
  \draw[assoc] (ua) -- (uc08.west);
  \draw[assoc] (ua) -- (uc09.west);
\end{tikzpicture}}
\caption{Generalização dos atores autenticados e casos de uso comuns}
\label{fig:uc-geral}
\end{figure}

\begin{figure}[H]
\centering
\resizebox{0.85\textwidth}{!}{%
\begin{tikzpicture}
  \ator{vis}{0,2.4}{Visitante}
  \draw[sistema] (2.4,-2.4) rectangle (11.4,7.4);
  \node[font=\small\bfseries] at (6.9,7.0) {Sistema de Gestão de Festivais};
  \node[uc] (uc01) at (5.6,6.0)  {UC01 Consultar lista de festivais};
  \node[uc] (uc02) at (5.6,4.5)  {UC02 Pesquisar festivais pelo nome};
  \node[uc] (uc03) at (5.6,3.0)  {UC03 Filtrar por data ou localização};
  \node[uc] (uc04) at (5.6,1.5)  {UC04 Consultar detalhes do festival};
  \node[uc] (uc05) at (5.6,0.0)  {UC05 Consultar programação,\\palcos e artistas};
  \node[uc] (uc06) at (5.6,-1.5) {UC06 Registar conta};
  \draw[assoc] (vis) -- (uc01.west);
  \draw[assoc] (vis) -- (uc04.west);
  \draw[assoc] (vis) -- (uc06.west);
  \draw[relacao] (uc02.north) -- node[relnome]{\extd} (uc01.south);
  \draw[relacao] (uc03.east) to[out=0,in=0,looseness=1.4]
        node[relnome]{\extd} (uc01.east);
  \draw[relacao] (uc05.north) -- node[relnome]{\extd} (uc04.south);
\end{tikzpicture}}
\caption{Casos de uso do Visitante}
\label{fig:uc-visitante}
\end{figure}

\begin{figure}[H]
\centering
\resizebox{0.85\textwidth}{!}{%
\begin{tikzpicture}
  \ator{par}{0,2.4}{Participante}
  \draw[sistema] (2.4,-0.9) rectangle (11.4,7.4);
  \node[font=\small\bfseries] at (6.9,7.0) {Sistema de Gestão de Festivais};
  \node[uc] (uc10) at (5.6,6.0) {UC10 Inscrever-se num festival};
  \node[uc] (uc11) at (5.6,4.5) {UC11 Efetuar pagamento simulado};
  \node[uc] (uc12) at (5.6,3.0) {UC12 Consultar as minhas inscrições};
  \node[uc] (uc13) at (5.6,1.5) {UC13 Consultar bilhete QR};
  \node[uc] (uc14) at (5.6,0.0) {UC14 Cancelar inscrição};
  \draw[assoc] (par) -- (uc10.west);
  \draw[assoc] (par) -- (uc12.west);
  \draw[relacao] (uc11.north) -- node[relnome]{\extd} (uc10.south);
  \draw[relacao] (uc11.south) -- node[relnome]{\extd} (uc12.north);
  \draw[relacao] (uc13.north) -- node[relnome]{\extd} (uc12.south);
  \draw[relacao] (uc14.east) to[out=0,in=0,looseness=1.4]
        node[relnome]{\extd} (uc12.east);
\end{tikzpicture}}
\caption{Casos de uso do Participante}
\label{fig:uc-participante}
\end{figure}

\begin{figure}[H]
\centering
\resizebox{0.85\textwidth}{!}{%
\begin{tikzpicture}
  \ator{org}{0,3.5}{Organizador}
  \draw[sistema] (2.4,-0.8) rectangle (11.4,8.6);
  \node[font=\small\bfseries] at (6.9,8.2) {Sistema de Gestão de Festivais};
  \node[uc] (uc15) at (6.9,7.2) {UC15 Criar festival};
  \node[uc] (uc16) at (6.9,6.0) {UC16 Editar festival e alterar estado};
  \node[uc] (uc17) at (6.9,4.8) {UC17 Gerir palcos};
  \node[uc] (uc18) at (6.9,3.6) {UC18 Gerir artistas do festival};
  \node[uc] (uc19) at (6.9,2.4) {UC19 Gerir programação};
  \node[uc] (uc20) at (6.9,1.2) {UC20 Gerir Staff do festival};
  \node[uc] (uc21) at (6.9,0.0) {UC21 Consultar inscrições do festival};
  \foreach \x in {uc15,uc16,uc17,uc18,uc19,uc20,uc21}
    \draw[assoc] (org) -- (\x.west);
\end{tikzpicture}}
\caption{Casos de uso do Organizador}
\label{fig:uc-organizador}
\end{figure}

\begin{figure}[H]
\centering
\resizebox{0.85\textwidth}{!}{%
\begin{tikzpicture}
  \ator{stf}{0,2.4}{Staff}
  \draw[sistema] (2.4,-0.9) rectangle (11.4,7.4);
  \node[font=\small\bfseries] at (6.9,7.0) {Sistema de Gestão de Festivais};
  \node[uc] (uc22) at (5.6,6.0) {UC22 Consultar festivais associados};
  \node[uc] (uc23) at (5.6,4.5) {UC23 Consultar participantes inscritos};
  \node[uc] (uc25) at (5.6,3.0) {UC25 Validar entrada};
  \node[uc] (uc24) at (5.6,1.5) {UC24 Verificar bilhete (QR ou código)};
  \node[uc] (uc26) at (5.6,0.0) {UC26 Consultar histórico de entradas};
  \foreach \x in {uc22,uc23,uc25,uc24,uc26}
    \draw[assoc] (stf) -- (\x.west);
  \draw[relacao] (uc25.south) -- node[relnome]{\incl} (uc24.north);
\end{tikzpicture}}
\caption{Casos de uso do Staff}
\label{fig:uc-staff}
\end{figure}

\begin{figure}[H]
\centering
\resizebox{0.85\textwidth}{!}{%
\begin{tikzpicture}
  \ator{adm}{0,2.4}{Administrador}
  \draw[sistema] (2.4,-0.9) rectangle (11.4,7.4);
  \node[font=\small\bfseries] at (6.9,7.0) {Sistema de Gestão de Festivais};
  \node[uc] (uc27) at (5.6,6.0) {UC27 Gerir contas de utilizador};
  \node[uc] (uc28) at (5.6,4.5) {UC28 Atribuir perfil};
  \node[uc] (uc29) at (5.6,3.0) {UC29 Consultar todos os festivais};
  \node[uc] (uc30) at (5.6,1.5) {UC30 Ativar ou desativar festival};
  \node[uc] (uc31) at (5.6,0.0) {UC31 Consultar indicadores globais};
  \draw[assoc] (adm) -- (uc27.west);
  \draw[assoc] (adm) -- (uc29.west);
  \draw[assoc] (adm) -- (uc31.west);
  \draw[relacao] (uc28.north) -- node[relnome]{\extd} (uc27.south);
  \draw[relacao] (uc30.north) -- node[relnome]{\extd} (uc29.south);
\end{tikzpicture}}
\caption{Casos de uso do Administrador}
\label{fig:uc-admin}
\end{figure}

% =====================================================================
\section{Descrição de Cenários}

Cada caso de uso é descrito pela pré-condição, pelo fluxo principal, pela
pós-condição e pelos fluxos alternativos. As referências RNxx remetem para as
regras de negócio da Secção~\ref{sec:regras}.

\subsection{Visitante}

\cenario{UC01}{Consultar lista de festivais}{Visitante}
{Nenhuma.}
{1.~O visitante acede à página inicial ou ao menu \emph{Festivais}.\newline
 2.~O sistema apresenta, com paginação, os festivais publicados e ativos
 (RN03), com nome, localização, datas e preço.}
{Lista de festivais apresentada.}
{Não existem festivais disponíveis: o sistema informa o visitante.}

\cenario{UC02}{Pesquisar festivais pelo nome}{Visitante}
{Lista de festivais apresentada (UC01).}
{1.~O visitante escreve parte do nome no campo de pesquisa.\newline
 2.~O sistema atualiza a lista por AJAX à medida que o texto é escrito.}
{Lista restrita aos festivais cujo nome contém o texto.}
{Nenhum festival corresponde à pesquisa: o sistema mostra uma mensagem.}

\cenario{UC03}{Filtrar festivais por data ou localização}{Visitante}
{Lista de festivais apresentada (UC01).}
{1.~O visitante indica um intervalo de datas e/ou uma localização.\newline
 2.~O sistema atualiza a lista com os festivais que decorrem nesse intervalo
 e/ou local.}
{Lista filtrada apresentada.}
{Data final anterior à inicial: o sistema indica o erro.
 Nenhum resultado: o sistema mostra uma mensagem.}

\cenario{UC04}{Consultar detalhes de um festival}{Visitante}
{O festival está publicado e ativo.}
{1.~O visitante seleciona um festival na lista.\newline
 2.~O sistema apresenta a descrição, localização, datas, preço, período de
 inscrições, lugares disponíveis e a opção \emph{Inscrever}.}
{Detalhes do festival apresentados.}
{O festival deixou de estar disponível: o sistema informa e volta à lista.}

\cenario{UC05}{Consultar programação, palcos e artistas}{Visitante}
{Detalhes do festival apresentados (UC04).}
{1.~O visitante abre o separador \emph{Programação}.\newline
 2.~O sistema apresenta as atuações agrupadas por dia e por palco, com
 artista e horário.\newline
 3.~O visitante pode selecionar um artista para ver o género e a biografia.}
{Programação apresentada.}
{Programação ainda não definida: o sistema mostra ``Programação em breve''.}

\cenario{UC06}{Registar conta}{Visitante}
{O visitante não tem sessão iniciada.}
{1.~O visitante preenche nome, email, palavra-passe e confirmação.\newline
 2.~O sistema valida os dados e a unicidade do email.\newline
 3.~O sistema cria a conta com perfil Participante (RN01), guarda apenas o
 \emph{hash} da palavra-passe e inicia sessão.}
{Conta ativa com perfil Participante; sessão iniciada.}
{Email já registado, dados inválidos ou palavras-passe diferentes: o sistema
 assinala os campos com erro e não cria a conta.}

\subsection{Utilizador Autenticado}

\cenario{UC07}{Iniciar sessão}{Utilizador Autenticado (qualquer perfil)}
{O utilizador tem uma conta registada.}
{1.~O utilizador introduz email e palavra-passe.\newline
 2.~O sistema verifica a palavra-passe contra o \emph{hash} guardado e o
 estado da conta.\newline
 3.~O sistema cria a sessão e redireciona para a área do respetivo perfil.}
{Sessão iniciada com as permissões do perfil.}
{Credenciais inválidas: mensagem genérica, sem indicar qual o campo errado.
 Conta desativada (RN02): acesso recusado.}

\cenario{UC08}{Terminar sessão}{Utilizador Autenticado (qualquer perfil)}
{Sessão iniciada.}
{1.~O utilizador seleciona \emph{Sair}.\newline
 2.~O sistema destrói a sessão e apresenta a página inicial.}
{Sessão terminada.}
{A sessão já tinha expirado: o sistema apresenta a página inicial.}

\cenario{UC09}{Editar perfil}{Utilizador Autenticado (qualquer perfil)}
{Sessão iniciada.}
{1.~O utilizador abre \emph{O meu perfil}.\newline
 2.~Altera o nome, o idioma preferido ou a palavra-passe (indicando a
 atual).\newline
 3.~O sistema valida e guarda as alterações.}
{Dados atualizados; a interface passa a usar o idioma escolhido.}
{Palavra-passe atual incorreta ou nova palavra-passe inválida: alteração
 recusada.}

\subsection{Participante}

\cenario{UC10}{Inscrever-se num festival}{Participante}
{Participante autenticado; festival publicado, ativo, com inscrições abertas e
 lugares disponíveis (RN05).}
{1.~No detalhe do festival, o participante seleciona \emph{Inscrever}.\newline
 2.~O sistema apresenta o resumo e o valor a pagar.\newline
 3.~O participante confirma.\newline
 4.~O sistema cria a inscrição no estado \emph{pendente}, com o valor igual ao
 preço atual (RN07), e propõe o pagamento (UC11).}
{Inscrição pendente de pagamento.}
{Inscrições fechadas ou lotação esgotada: inscrição recusada.
 Já existe inscrição pendente ou confirmada (RN06): o sistema mostra a
 inscrição existente.
 Inscrição anterior cancelada: é criada uma nova inscrição e a anterior
 fica no histórico.}

\cenario{UC11}{Efetuar pagamento simulado}{Participante}
{Existe uma inscrição pendente do participante.}
{1.~O participante escolhe o método (cartão, MB WAY ou referência
 Multibanco, todos simulados) e confirma.\newline
 2.~O sistema regista a tentativa de pagamento e o resultado.\newline
 3.~Se o pagamento é aprovado, o sistema, numa única transação, confirma a
 inscrição e emite o bilhete com código QR único (RN08, RN09).}
{Pagamento aprovado, inscrição confirmada e bilhete emitido.}
{Pagamento recusado: a tentativa fica registada, a inscrição continua
 pendente e o participante pode tentar de novo.
 A lotação esgotou entretanto: o pagamento não é processado e o participante
 é informado.}

\cenario{UC12}{Consultar as minhas inscrições}{Participante}
{Participante autenticado.}
{1.~O participante abre \emph{As minhas inscrições}.\newline
 2.~O sistema lista as inscrições com festival, datas, valor e estado.\newline
 3.~Em cada inscrição, o participante pode pagar (UC11), ver o bilhete (UC13)
 ou cancelar (UC14), conforme o estado.}
{Lista de inscrições apresentada.}
{Não existem inscrições: o sistema sugere a lista de festivais.}

\cenario{UC13}{Consultar bilhete QR}{Participante}
{Inscrição confirmada com bilhete não revogado.}
{1.~O participante seleciona \emph{Ver bilhete}.\newline
 2.~O sistema apresenta o código QR, o código manual, o nome do participante
 e os dados do festival, com opção de impressão.}
{Bilhete apresentado.}
{Inscrição pendente ou cancelada: o bilhete não está disponível.}

\cenario{UC14}{Cancelar inscrição}{Participante}
{Inscrição pendente ou confirmada, dentro do prazo de cancelamento e sem
 entrada registada (RN10).}
{1.~O participante seleciona \emph{Cancelar} e o sistema pede
 confirmação.\newline
 2.~O participante confirma.\newline
 3.~O sistema, numa transação, passa a inscrição a \emph{cancelada}, regista a
 data e revoga o bilhete, se existir.}
{Inscrição cancelada, bilhete revogado e lugar libertado.}
{Prazo ultrapassado ou entrada já registada: cancelamento recusado.
 O participante desiste na confirmação: nada é alterado.}

\subsection{Organizador}

\cenario{UC15}{Criar festival}{Organizador}
{Sessão iniciada com perfil Organizador.}
{1.~O organizador seleciona \emph{Novo festival}.\newline
 2.~Introduz nome, descrição, localização, datas de início e fim, período
 de inscrições, data limite de cancelamento, capacidade e preço.\newline
 3.~O sistema valida os dados e cria o festival no estado \emph{rascunho},
 tendo o organizador como responsável.}
{Festival criado em rascunho, ainda não visível ao público.}
{Dados em falta ou incoerentes (fim antes do início, fecho das inscrições
 depois do fim do festival, capacidade ou preço negativos): o sistema
 assinala os erros.}

\cenario{UC16}{Editar festival e alterar estado}{Organizador}
{O organizador é o responsável pelo festival (RN04).}
{1.~O organizador abre o festival e altera os dados, ou\newline
 2.~altera o estado: publicar (rascunho para publicado), cancelar ou
 terminar.\newline
 3.~O sistema valida e guarda.}
{Dados ou estado do festival atualizados.}
{Capacidade inferior ao número de inscrições confirmadas: alteração recusada.
 Festival cancelado ou terminado: não pode ser editado.
 Festival desativado pelo Administrador: pode ser publicado mas continua
 invisível até ser reativado (RN03).}

\cenario{UC17}{Gerir palcos}{Organizador}
{O organizador é o responsável pelo festival.}
{1.~O organizador abre a lista de palcos do festival.\newline
 2.~Adiciona, renomeia ou remove um palco.\newline
 3.~O sistema valida e guarda.}
{Palcos do festival atualizados.}
{Nome repetido no mesmo festival: recusado.
 Palco com atuações marcadas (RN14): remoção recusada.}

\cenario{UC18}{Gerir artistas do festival}{Organizador}
{O organizador é o responsável pelo festival.}
{1.~O organizador pesquisa um artista no catálogo ou cria um novo (nome,
 género e biografia).\newline
 2.~Associa o artista ao festival ou remove a associação.\newline
 3.~O sistema guarda as alterações.}
{Lista de artistas do festival atualizada.}
{Artista já associado: o sistema informa.
 Artista com atuações marcadas no festival (RN14): remoção recusada.}

\cenario{UC19}{Gerir programação}{Organizador}
{O organizador é o responsável pelo festival, que tem pelo menos um palco e um
 artista associado.}
{1.~O organizador seleciona \emph{Nova atuação}.\newline
 2.~Escolhe o artista, o palco, a data e hora de início e de fim.\newline
 3.~O sistema valida a atuação (RN13) e guarda. As atuações existentes podem
 ser editadas ou removidas.}
{Programação atualizada; visível ao público se o festival estiver publicado.}
{Sobreposição com outra atuação no mesmo palco, horário fora das datas do
 festival ou fim anterior ao início: o sistema recusa e indica o motivo.}

\cenario{UC20}{Gerir Staff do festival}{Organizador}
{O organizador é o responsável pelo festival.}
{1.~O organizador pesquisa utilizadores com perfil Staff por nome ou
 email.\newline
 2.~Associa o utilizador ao festival, ou remove uma associação existente.\newline
 3.~O sistema guarda; a remoção desativa a associação (RN12).}
{Lista de Staff do festival atualizada.}
{Utilizador inexistente ou sem perfil Staff: recusado.
 Associação já ativa: o sistema informa.
 Associação inativa: é reativada.}

\cenario{UC21}{Consultar inscrições do festival}{Organizador}
{O organizador é o responsável pelo festival.}
{1.~O organizador abre o separador \emph{Inscrições}.\newline
 2.~O sistema lista os participantes, o estado e o valor de cada inscrição,
 com filtro por estado.\newline
 3.~O sistema apresenta indicadores: inscrições confirmadas, ocupação face à
 capacidade, receita simulada e gráfico de inscrições por dia.}
{Inscrições e indicadores apresentados.}
{Não existem inscrições: o sistema informa.}

\subsection{Staff}

\cenario{UC22}{Consultar festivais associados}{Staff}
{Sessão iniciada com perfil Staff.}
{1.~O sistema lista os festivais com associação ativa, destacando os que
 decorrem no próprio dia.\newline
 2.~O membro do Staff seleciona o festival em que vai trabalhar.}
{Festival selecionado para as operações seguintes.}
{Não existem associações ativas: o sistema informa.}

\cenario{UC23}{Consultar participantes inscritos}{Staff}
{Membro do Staff associado ao festival (RN12).}
{1.~O sistema lista os participantes com inscrição confirmada.\newline
 2.~O membro do Staff pesquisa por nome, email ou código do bilhete.\newline
 3.~Para cada participante, o sistema indica se a entrada já foi registada.}
{Lista de participantes apresentada.}
{Participante não encontrado: o sistema informa.}

\cenario{UC24}{Verificar bilhete}{Staff}
{Membro do Staff associado ao festival selecionado.}
{1.~O membro do Staff lê o código QR com a câmara ou introduz o código
 manual.\newline
 2.~O sistema verifica o bilhete por AJAX (RN11) e apresenta o resultado
 (válido ou inválido, com o motivo) e o nome do participante.}
{Resultado apresentado; nada é registado.}
{Código inexistente, bilhete de outro festival, bilhete revogado, festival
 que não decorre nesse dia ou bilhete já utilizado (com data e hora da
 entrada anterior): o bilhete é indicado como inválido.}

\cenario{UC25}{Validar entrada}{Staff}
{Membro do Staff associado ao festival selecionado.}
{1.~O sistema verifica o bilhete (\incl{} UC24).\newline
 2.~Se o bilhete é válido, o membro do Staff confirma a entrada.\newline
 3.~O sistema regista a entrada (bilhete, membro do Staff, data e hora).}
{Entrada registada; o bilhete fica utilizado.}
{Bilhete inválido: entrada recusada com o motivo.
 Leitura simultânea do mesmo bilhete noutro posto: a restrição de unicidade
 da base de dados rejeita o segundo registo e o sistema indica ``bilhete já
 utilizado''.}

\cenario{UC26}{Consultar histórico de entradas}{Staff}
{Membro do Staff associado ao festival.}
{1.~O sistema lista as entradas do festival com data e hora, participante e
 membro do Staff que validou.\newline
 2.~O sistema mostra o total de entradas face ao total de inscrições
 confirmadas.}
{Histórico apresentado.}
{Não existem entradas registadas: o sistema informa.}

\subsection{Administrador}

\cenario{UC27}{Gerir contas de utilizador}{Administrador}
{Sessão iniciada com perfil Administrador.}
{1.~O sistema lista os utilizadores, com pesquisa e filtros por perfil e
 estado.\newline
 2.~O administrador cria uma conta, edita dados, ativa ou desativa uma
 conta.\newline
 3.~O sistema valida e guarda.}
{Contas atualizadas.}
{Email já usado por outra conta: recusado.
 O administrador tenta desativar a própria conta: recusado.}

\cenario{UC28}{Atribuir perfil}{Administrador}
{Conta selecionada em UC27.}
{1.~O administrador escolhe o novo perfil e confirma.\newline
 2.~O sistema atualiza o perfil da conta.}
{Perfil alterado; as novas permissões aplicam-se na operação seguinte do
 utilizador.}
{Retirar o perfil Organizador a quem tem festivais publicados, ou retirar o
 perfil Administrador ao último administrador: recusado.
 Retirar o perfil Staff: as associações a festivais são desativadas.}

\cenario{UC29}{Consultar todos os festivais}{Administrador}
{Sessão iniciada com perfil Administrador.}
{1.~O sistema lista todos os festivais, em qualquer estado, com organizador,
 estado, indicador ativo e número de inscrições.\newline
 2.~O administrador pode filtrar por estado ou organizador.}
{Lista apresentada.}
{Não existem festivais: o sistema informa.}

\cenario{UC30}{Ativar ou desativar festival}{Administrador}
{Festival selecionado em UC29.}
{1.~O administrador seleciona \emph{Desativar} (ou \emph{Ativar}) e
 confirma.\newline
 2.~O sistema atualiza o indicador ativo do festival.}
{Um festival desativado deixa de ser visível e de aceitar inscrições (RN03);
 as inscrições existentes mantêm-se.}
{O administrador desiste na confirmação: nada é alterado.}

\cenario{UC31}{Consultar indicadores globais}{Administrador}
{Sessão iniciada com perfil Administrador.}
{1.~O administrador abre o \emph{dashboard} e escolhe um período.\newline
 2.~O sistema apresenta, com gráficos: utilizadores por perfil, festivais por
 estado, inscrições por estado, pagamentos aprovados e recusados e receita
 simulada.}
{Indicadores apresentados.}
{Não existem dados no período: o sistema informa.}

% =====================================================================
\section{Diagramas de Estados}

Os diagramas de estados das Figuras~\ref{fig:estados-inscricao}
e~\ref{fig:estados-festival} resumem o ciclo de vida das duas entidades
centrais e correspondem às colunas \texttt{estado} das tabelas
\texttt{inscricao} e \texttt{festival}.

\begin{figure}[H]
\centering
\resizebox{0.95\textwidth}{!}{%
\begin{tikzpicture}
  \fill (0,0) circle (0.13);
  \node[estado] (pen) at (3,0)  {pendente};
  \node[estado] (con) at (10,0) {confirmada};
  \node[estado] (can) at (6.5,-3.2) {cancelada};
  \draw[trans] (0.13,0) -- node[translab,above]{inscrever} (pen);
  \draw[trans] (pen) -- node[translab,above]{pagamento aprovado\\/ emite bilhete} (con);
  \draw[trans] (pen.north west) .. controls +(-0.4,1.4) and +(0.4,1.4) ..
        node[translab,above]{pagamento recusado} (pen.north east);
  \draw[trans] (pen.south) -- node[translab,left]{cancelar} (can.west);
  \draw[trans] (con.south) -- node[translab,right]{cancelar [no prazo\\e sem entrada]\\/ revoga bilhete} (can.east);
  \draw[thick] (6.5,-5.2) circle (0.2);
  \fill (6.5,-5.2) circle (0.12);
  \draw[trans] (can.south) -- (6.5,-5.0);
\end{tikzpicture}}
\caption{Diagrama de estados da inscrição (uma nova inscrição após
cancelamento é uma nova instância)}
\label{fig:estados-inscricao}
\end{figure}

\begin{figure}[H]
\centering
\resizebox{0.9\textwidth}{!}{%
\begin{tikzpicture}
  \fill (0,0) circle (0.13);
  \node[estado] (ras) at (3,0)   {rascunho};
  \node[estado] (pub) at (8,0)   {publicado};
  \node[estado] (ter) at (12.5,0) {terminado};
  \node[estado] (can) at (5.5,-3) {cancelado};
  \draw[trans] (0.13,0) -- node[translab,above]{criar} (ras);
  \draw[trans] (ras) -- node[translab,above]{publicar} (pub);
  \draw[trans] (pub) -- node[translab,above]{terminar} (ter);
  \draw[trans] (ras.south) -- node[translab,left]{cancelar} (can.west);
  \draw[trans] (pub.south) -- node[translab,right]{cancelar} (can.east);
\end{tikzpicture}}
\caption{Diagrama de estados do festival (independente do indicador
\emph{ativo}, gerido pelo Administrador)}
\label{fig:estados-festival}
\end{figure}

\chapter{Base de Dados}
\label{cap:bd}

% ---------- Macros locais do diagrama ER ----------
% \entidade{nome-nó}{posição do topo}{título}{linhas: chave & coluna & tipo \\}
\newcommand{\entidade}[4]{%
  \node[draw, thick, fill=white, inner sep=0pt, anchor=north, font=\scriptsize]
    (#1) at (#2)
    {\begin{tabular}{@{\,}l@{\;\;}l@{\;\;}l@{\,}}
       \multicolumn{3}{c}{\rule{0pt}{2.6ex}\small\textbf{#3}}\\[0.4ex] \hline
       #4
     \end{tabular}};}
% \relER{origem}{destino}{nome}{cardinalidade origem}{cardinalidade destino}
\newcommand{\relER}[5]{%
  \draw[thin] (#1) --
    node[pos=0.5, fill=white, inner sep=1pt, font=\scriptsize\itshape]{#3}
    node[pos=0.1, auto, font=\scriptsize\bfseries]{#4}
    node[pos=0.9, auto, font=\scriptsize\bfseries]{#5} (#2);}

\section{Modelo Entidade--Relação}

O modelo de dados tem doze tabelas. As decisões principais são:
\begin{itemize}
  \item \textbf{perfil} é uma tabela (e não um tipo enumerado) para permitir
        descrições e acrescentar perfis sem alterar a estrutura;
  \item \textbf{artista} é um catálogo partilhado; a ligação aos festivais é
        de muitos-para-muitos através de \textbf{festival\_artista};
  \item uma \textbf{atuacao} referencia o par festival--artista e o par
        palco--festival, o que garante ao nível da base de dados que o artista
        e o palco pertencem ao mesmo festival (RN13);
  \item \textbf{pagamento} é separado da inscrição para registar todas as
        tentativas (aprovadas e recusadas), mas só uma pode ser aprovada
        (RN08);
  \item \textbf{bilhete} é separado da inscrição porque só existe depois da
        aprovação do pagamento e pode ser revogado (RN09, RN10);
  \item \textbf{entrada} tem no máximo um registo por bilhete, o que impede a
        reutilização (RN11);
  \item a associação do Staff (\textbf{festival\_staff}) é desativada e não
        apagada, para manter o histórico das entradas validadas (RN12).
\end{itemize}

\begin{figure}[H]
\centering
\resizebox{\textwidth}{!}{%
\begin{tikzpicture}
  % ---- Linha A ----
  \entidade{perfil}{0,0}{perfil}{%
    PK & codigo & VARCHAR(20)\\
       & descricao & VARCHAR(100)\\}
  \entidade{utilizador}{6.5,0}{utilizador}{%
    PK & id\_utilizador & BIGINT\\
       & nome & VARCHAR(120)\\
    UK & email & VARCHAR(254)\\
       & password\_hash & VARCHAR(255)\\
    FK & perfil\_codigo & VARCHAR(20)\\
       & idioma & CHAR(2)\\
       & ativo & BOOLEAN\\
       & criado\_em & DATETIME\\}
  \entidade{fstaff}{13.5,0}{festival\_staff}{%
    PK,FK & id\_festival & BIGINT\\
    PK,FK & id\_utilizador & BIGINT\\
          & associado\_em & DATETIME\\
          & ativo & BOOLEAN\\}
  % ---- Linha B ----
  \entidade{inscricao}{0,-5.5}{inscricao}{%
    PK & id\_inscricao & BIGINT\\
    FK & id\_participante & BIGINT\\
    FK & id\_festival & BIGINT\\
       & estado & ENUM\\
       & valor & DECIMAL(10,2)\\
       & criada\_em & DATETIME\\
       & confirmada\_em & DATETIME\\
       & cancelada\_em & DATETIME\\
    UK & id\_festival\_ativa & BIGINT (gerada)\\}
  \entidade{festival}{6.5,-5.5}{festival}{%
    PK & id\_festival & BIGINT\\
    FK & id\_organizador & BIGINT\\
       & nome & VARCHAR(160)\\
       & descricao & TEXT\\
       & localizacao & VARCHAR(200)\\
       & data\_inicio & DATE\\
       & data\_fim & DATE\\
       & inscricoes\_abrem & DATETIME\\
       & inscricoes\_fecham & DATETIME\\
       & cancelamento\_ate & DATETIME\\
       & capacidade & INT\\
       & preco & DECIMAL(10,2)\\
       & estado & ENUM\\
       & ativo & BOOLEAN\\
       & criado\_em & DATETIME\\}
  \entidade{fartista}{13.5,-5.5}{festival\_artista}{%
    PK,FK & id\_festival & BIGINT\\
    PK,FK & id\_artista & BIGINT\\}
  \entidade{artista}{19.5,-5.5}{artista}{%
    PK & id\_artista & BIGINT\\
       & nome & VARCHAR(160)\\
       & genero & VARCHAR(80)\\
       & biografia & TEXT\\}
  % ---- Linha C ----
  \entidade{pagamento}{-3,-14}{pagamento}{%
    PK & id\_pagamento & BIGINT\\
    FK & id\_inscricao & BIGINT\\
       & metodo & ENUM\\
       & resultado & ENUM\\
       & valor & DECIMAL(10,2)\\
       & criado\_em & DATETIME\\
    UK & inscricao\_aprovada & BIGINT (gerada)\\}
  \entidade{bilhete}{3,-14}{bilhete}{%
    PK & id\_bilhete & BIGINT\\
    FK,UK & id\_inscricao & BIGINT\\
    UK & codigo\_qr & CHAR(64)\\
       & emitido\_em & DATETIME\\
       & revogado\_em & DATETIME\\}
  \entidade{palco}{9.5,-14}{palco}{%
    PK & id\_palco & BIGINT\\
    FK & id\_festival & BIGINT\\
       & nome & VARCHAR(100)\\}
  \entidade{atuacao}{16,-14}{atuacao}{%
    PK & id\_atuacao & BIGINT\\
    FK & id\_festival & BIGINT\\
    FK & id\_palco & BIGINT\\
    FK & id\_artista & BIGINT\\
       & inicio & DATETIME\\
       & fim & DATETIME\\}
  % ---- Linha D ----
  \entidade{entrada}{3,-19}{entrada}{%
    PK & id\_entrada & BIGINT\\
    FK,UK & id\_bilhete & BIGINT\\
    FK & id\_staff & BIGINT\\
       & validada\_em & DATETIME\\}

  % ---- Relações ----
  \relER{perfil.east}{utilizador.west}{classifica}{1}{N}
  \relER{utilizador.east}{fstaff.west}{trabalha\_em}{1}{N}
  \relER{fstaff.south}{festival.north east}{tem\_staff}{N}{1}
  \relER{utilizador.south}{festival.north}{organiza}{1}{N}
  \relER{utilizador.south west}{inscricao.north east}{faz}{1}{N}
  \relER{festival.west}{inscricao.east}{recebe}{1}{N}
  \relER{festival.east}{fartista.west}{associa}{1}{N}
  \relER{artista.west}{fartista.east}{participa}{1}{N}
  \relER{festival.south}{palco.north}{tem}{1}{N}
  \relER{palco.east}{atuacao.west}{acolhe}{1}{N}
  \relER{fartista.south}{atuacao.north}{programa}{1}{N}
  \relER{inscricao.south west}{pagamento.north}{tem}{1}{N}
  \relER{inscricao.south east}{bilhete.north}{emite}{1}{0..1}
  \relER{bilhete.south}{entrada.north}{origina}{1}{0..1}
  % utilizador (Staff) valida entrada: percurso pela esquerda
  \coordinate (cantoValida) at ($(entrada.west)+(-7.8,0)$);
  \draw[thin] (utilizador.north) -- ++(0,1) -| (cantoValida) -- (entrada.west);
  \node[fill=white, inner sep=1pt, font=\scriptsize\itshape, rotate=90]
    at ($(cantoValida)+(0,11)$) {valida};
  \node[font=\scriptsize\bfseries] at ($(utilizador.north)+(0.25,0.35)$) {1};
  \node[font=\scriptsize\bfseries] at ($(entrada.west)+(-0.35,0.25)$) {N};
\end{tikzpicture}}
\caption{Modelo entidade--relação (PK: chave primária; FK: chave estrangeira;
UK: valor único)}
\label{fig:er}
\end{figure}

\section{Modelo Relacional}

No esquema seguinte, as chaves primárias estão \underline{sublinhadas} e as
chaves estrangeiras estão em \textit{itálico}, com a tabela referenciada
indicada por $\rightarrow$.

\begin{itemize}\small
  \item \textbf{perfil}(\underline{codigo}, descricao)
  \item \textbf{utilizador}(\underline{id\_utilizador}, nome, email, password\_hash,
        \textit{perfil\_codigo}$\rightarrow$perfil, idioma, ativo, criado\_em)
  \item \textbf{festival}(\underline{id\_festival}, \textit{id\_organizador}$\rightarrow$utilizador,
        nome, descricao, localizacao, data\_inicio, data\_fim, inscricoes\_abrem,
        inscricoes\_fecham, cancelamento\_ate, capacidade, preco, estado, ativo,
        criado\_em)
  \item \textbf{festival\_staff}(\underline{\textit{id\_festival}}$\rightarrow$festival,
        \underline{\textit{id\_utilizador}}$\rightarrow$utilizador, associado\_em, ativo)
  \item \textbf{artista}(\underline{id\_artista}, nome, genero, biografia)
  \item \textbf{festival\_artista}(\underline{\textit{id\_festival}}$\rightarrow$festival,
        \underline{\textit{id\_artista}}$\rightarrow$artista)
  \item \textbf{palco}(\underline{id\_palco}, \textit{id\_festival}$\rightarrow$festival, nome)
  \item \textbf{atuacao}(\underline{id\_atuacao}, \textit{id\_festival}, \textit{id\_palco},
        \textit{id\_artista}, inicio, fim), com
        (\textit{id\_festival}, \textit{id\_artista})$\rightarrow$festival\_artista e
        (\textit{id\_palco}, \textit{id\_festival})$\rightarrow$palco
  \item \textbf{inscricao}(\underline{id\_inscricao}, \textit{id\_participante}$\rightarrow$utilizador,
        \textit{id\_festival}$\rightarrow$festival, estado, valor, criada\_em,
        confirmada\_em, cancelada\_em, id\_festival\_ativa)
  \item \textbf{pagamento}(\underline{id\_pagamento}, \textit{id\_inscricao}$\rightarrow$inscricao,
        metodo, resultado, valor, criado\_em, inscricao\_aprovada)
  \item \textbf{bilhete}(\underline{id\_bilhete}, \textit{id\_inscricao}$\rightarrow$inscricao,
        codigo\_qr, emitido\_em, revogado\_em)
  \item \textbf{entrada}(\underline{id\_entrada}, \textit{id\_bilhete}$\rightarrow$bilhete,
        \textit{id\_staff}$\rightarrow$utilizador, validada\_em)
\end{itemize}

\noindent Domínios dos tipos enumerados:
\begin{itemize}\small
  \item \texttt{perfil.codigo}: PARTICIPANTE, ORGANIZADOR, STAFF, ADMIN;
  \item \texttt{festival.estado}: rascunho, publicado, cancelado, terminado;
  \item \texttt{inscricao.estado}: pendente, confirmada, cancelada;
  \item \texttt{pagamento.metodo}: cartao, mbway, multibanco;
        \texttt{pagamento.resultado}: aprovado, recusado;
  \item \texttt{utilizador.idioma}: pt, en.
\end{itemize}

\section{Restrições de Integridade}

A Tabela~\ref{tab:integridade} relaciona as restrições da base de dados com
as regras de negócio que garantem. As regras que não se podem exprimir apenas
com chaves e restrições (por exemplo, a sobreposição de atuações) são
verificadas na aplicação, dentro de uma transação, ou por \emph{trigger}.

\begin{longtable}{|p{6.0cm}|p{6.0cm}|l|}
\caption{Restrições de integridade}\label{tab:integridade}\\
\hline
\textbf{Restrição} & \textbf{Objetivo} & \textbf{Regra} \\ \hline
\endfirsthead
\hline
\textbf{Restrição} & \textbf{Objetivo (cont.)} & \textbf{Regra} \\ \hline
\endhead
UNIQUE \texttt{utilizador(email)} & Uma conta por email. & RN01 \\ \hline
CHECK em \texttt{festival}: \texttt{data\_fim >= data\_inicio},
  \texttt{inscricoes\_fecham > inscricoes\_abrem}, \texttt{capacidade > 0},
  \texttt{preco >= 0} & Dados do festival coerentes. & -- \\ \hline
Coluna gerada \texttt{id\_festival\_ativa} = \texttt{id\_festival} se a
  inscrição não está cancelada, senão NULL; UNIQUE
  \texttt{(id\_participante, id\_festival\_ativa)} & No máximo uma inscrição
  não cancelada por participante e festival; as canceladas ficam como
  histórico. & RN06 \\ \hline
Coluna gerada \texttt{inscricao\_aprovada} = \texttt{id\_inscricao} se
  \texttt{resultado = 'aprovado'}, senão NULL; UNIQUE & No máximo um pagamento
  aprovado por inscrição (os NULL não colidem). & RN08 \\ \hline
UNIQUE \texttt{bilhete(id\_inscricao)} e \texttt{bilhete(codigo\_qr)} & Um
  bilhete por inscrição e código QR único. & RN09 \\ \hline
UNIQUE \texttt{entrada(id\_bilhete)} & Cada bilhete só dá uma entrada, mesmo
  com leituras simultâneas. & RN11 \\ \hline
UNIQUE \texttt{palco(id\_festival, nome)} e \texttt{palco(id\_palco,
  id\_festival)} & Nomes de palco únicos por festival; suporte à chave
  estrangeira composta. & RN13 \\ \hline
FK \texttt{atuacao(id\_festival, id\_artista)} $\rightarrow$
  \texttt{festival\_artista}; FK \texttt{atuacao(id\_palco, id\_festival)}
  $\rightarrow$ \texttt{palco}; CHECK \texttt{fim > inicio} & Artista e palco
  do mesmo festival; horário válido. & RN13 \\ \hline
ON DELETE RESTRICT nas chaves que referenciam \texttt{palco} e
  \texttt{festival\_artista} & Não remover palcos ou artistas com atuações.
  & RN14 \\ \hline
Verificação na aplicação ou por \emph{trigger} & Sem sobreposição de
  atuações no mesmo palco; atuação dentro das datas do festival. & RN13 \\ \hline
Verificação na aplicação & O Staff que valida está associado ao festival do
  bilhete; o organizador é o responsável pelo festival. & RN04, RN12 \\ \hline
\end{longtable}

\noindent Operações executadas em \textbf{transação}:
\begin{itemize}
  \item pagamento aprovado: inserir \texttt{pagamento}, confirmar a
        \texttt{inscricao} e inserir o \texttt{bilhete} (com nova verificação
        da lotação);
  \item cancelamento: atualizar a \texttt{inscricao} e revogar o
        \texttt{bilhete};
  \item validação de entrada: verificar o bilhete e inserir a
        \texttt{entrada}.
\end{itemize}

\noindent \textbf{Vistas} previstas: \texttt{vw\_festivais\_publicos},
\texttt{vw\_programacao}, \texttt{vw\_ocupacao\_festival} (confirmadas,
capacidade e entradas por festival) e \texttt{vw\_indicadores\_globais}.

\noindent O \texttt{codigo\_qr} é gerado com 32 \emph{bytes} aleatórios
criptograficamente seguros, representados em hexadecimal (64 caracteres), para
não ser previsível.

\section{Criação da Base de Dados e Dados de Exemplo}

\pendente{Descrever os scripts da pasta \texttt{bd/} (esquema, vistas e
procedimentos, dados de exemplo, testes) e referir o \emph{dump} entregue.}

\section{Queries SQL}

\pendente{Principais consultas por caso de uso (por exemplo: festivais
públicos com filtros (UC01--UC03), programação (UC05), inscrições do
participante (UC12), verificação de bilhete (UC24), ocupação e indicadores
(UC21, UC31)).}

\chapter{Design da Interface Gráfica}

\section{Design do Website}

\pendente{Estrutura geral das páginas, paleta de cores, tipografia, menu por
perfil e comportamento responsivo (Bootstrap). Indicar a seleção de idioma
PT/EN.}

\section{Formulários e Menus de Entrada de Dados (\emph{Wireframes})}

\pendente{\emph{Wireframes} dos principais formulários: registo e início de
sessão (UC06, UC07), inscrição e pagamento (UC10, UC11), criação de festival
(UC15), programação (UC19), validação de entrada (UC24, UC25) e gestão de
contas (UC27, UC28).}

\section{Formatos de Saída de Dados}

\pendente{Listagens (festivais, inscrições, participantes, entradas), bilhete
com código QR (UC13), programação por palco e dia (UC05) e \emph{dashboards}
com gráficos (UC21, UC31).}

\section{Fluxo entre Páginas (\emph{Wireflow})}

\pendente{Diagrama de navegação entre páginas para cada perfil, começando
pela página inicial do Visitante.}

\chapter{Implementação}

\section{Ferramentas e Tecnologias}

\pendente{Tabela com as ferramentas usadas e a respetiva função: PHP, MySQL 8
(PDO), HTML/CSS/Bootstrap, JavaScript/AJAX, biblioteca de leitura e geração
de QR, servidor local (XAMPP/WAMP/Docker), Git/GitHub, VSCode, Jira/Kanban.}

\section{Arquitetura (MVC)}

\pendente{Diagrama da arquitetura em camadas (\emph{View}, \emph{Controller},
\emph{Model}, base de dados) e estrutura de pastas do projeto.}

\section{Diagrama de Instalação}

\pendente{Cliente (navegador), servidor \emph{web} com PHP e servidor MySQL.}

\section{Autenticação, Controlo de Acesso e Multilinguismo}

\pendente{Sessões, \texttt{password\_hash}/\texttt{password\_verify},
verificação do perfil em cada \emph{controller}, e ficheiros de tradução PT/EN.}

\section{Funcionalidades Implementadas}

\pendente{Capturas de ecrã por perfil (Visitante, Participante, Organizador,
Staff, Administrador), indicando os casos de uso concretizados.}

\chapter{Testes e Validação}

\pendente{Tabela de casos de teste (ID, UC, passos, resultado esperado,
resultado obtido). Incluir, no mínimo, os casos críticos: inscrição duplicada
(RN06), pagamento recusado e repetido (RN08), lotação esgotada (RN05),
cancelamento fora do prazo (RN10), bilhete reutilizado (RN11), Staff não
associado (RN12) e sobreposição de atuações (RN13).}

\chapter{Conclusão}

\section{Sumário dos Resultados}

\pendente{Requisitos cumpridos e não cumpridos, face às Tabelas~\ref{tab:rf}
e~\ref{tab:rnf}.}

\section{Trabalho Futuro}

\pendente{Por exemplo: pagamentos reais, tipos de bilhete (diário, VIP),
aplicação móvel para o Staff, autenticação com Google/GitHub e notificações
por email.}


\appendix
\chapter{Anexos}

\section{Repositório GitHub}

\pendente{Ligação para o repositório e para o quadro Kanban/Jira.}

\section{Elementos Adicionais}

\pendente{Dump da base de dados, credenciais de demonstração por perfil e
instruções de instalação.}


\end{document}
